\noindent
\subsection{Dạng 1. Xác định cặp góc so le trong, cặp góc đồng vị. Tính số đo góc}

\begin{bt}
	Cho hình vẽ dưới đây, trong đó hai đường thẳng $m$ và $n$ bị cắt bởi
	đường thẳng $p$ lần lượt tại $S$ và $R$.
	
	\begin{enumerate}[label=\alph*)]
		\item Kể tên các cặp góc so le trong, các cặp góc đồng vị và các cặp
		góc trong cùng phía.
		
		\item Tính số đo các góc còn lại, biết
		$\widehat{R_4}=\widehat{S_2}=120^\circ$.
	\end{enumerate}

    \begin{center}
        \begin{tikzpicture}[
            scale=1.15,
            line cap=round,
            line join=round,
            font=\small,
            every node/.style={inner sep=1.2pt}
        ]
            %--------------------------------
            % Hai giao điểm
            %--------------------------------
            \coordinate (S) at (0,1.1);
            \coordinate (R) at (1.27,-1.1);

            %--------------------------------
            % Hai đường thẳng song song m, n
            %--------------------------------
            \draw[mauChuDao,line width=0.9pt]
                (-2.3,1.1)--(2.3,1.1)
                node[right] {$m$};

            \draw[mauChuDao,line width=0.9pt]
                (-0.8,-1.1)--(3.8,-1.1)
                node[right] {$n$};

            % Ký hiệu m // n nếu muốn
            % \node at (-1.6,0.85) {$\parallel$};

            %--------------------------------
            % Đường cắt p
            %--------------------------------
            \draw[mauChuDao,line width=0.9pt]
                ($(S)+(120:1.25)$)
                --
                ($(R)+(-60:1.25)$);

            \node[above left] at ($(S)+(120:1.20)$) {$p$};

            %--------------------------------
            % Hai giao điểm
            %--------------------------------
            \fill[red!75!black] (S) circle (1.5pt);
            \fill[red!75!black] (R) circle (1.5pt);

            \node[font=\bfseries]
                at ($(S)+(90:0.72)$) {$S$};

            \node[font=\bfseries]
                at ($(R)+(270:0.72)$) {$R$};

            %--------------------------------
            % Bốn góc tại S
            %--------------------------------
            \node at ($(S)+(150:0.50)$) {$1$};
            \node at ($(S)+(60:0.50)$)  {$2$};
            \node at ($(S)+(330:0.50)$) {$3$};
            \node at ($(S)+(240:0.50)$) {$4$};

            %--------------------------------
            % Bốn góc tại R
            %--------------------------------
            \node at ($(R)+(150:0.50)$) {$1$};
            \node at ($(R)+(60:0.50)$)  {$2$};
            \node at ($(R)+(330:0.50)$) {$3$};
            \node at ($(R)+(240:0.50)$) {$4$};

            %--------------------------------
            % Đánh dấu hai góc bằng nhau
            % Ví dụ: góc S_2 và R_4
            %--------------------------------
            \draw[mauChuDao,line width=1pt]
                ($(S)+(0:0.31)$)
                arc[start angle=0,end angle=120,radius=0.31];

            \draw[mauChuDao,line width=1pt]
                ($(R)+(180:0.31)$)
                arc[start angle=180,end angle=300,radius=0.31];

        \end{tikzpicture}
    \end{center}
\end{bt}

\begin{bt}
	Cho hình vẽ dưới đây, biết $a\parallel b$ và
	$\widehat{M_1}-\widehat{N_2}=50^\circ$. Tính $\widehat{M_2}$ và
	$\widehat{N_2}$.

	\begin{center}
		\begin{tikzpicture}[
			scale=1.15,
			line cap=round,
			line join=round,
			font=\small,
			every node/.style={inner sep=1.2pt}
		]
			\coordinate (M) at (0,-1);
			\coordinate (N) at (1.03,1.2);

			% Hai đường thẳng song song và cát tuyến c
			\draw[mauChuDao,line width=0.9pt]
				(-1.5,1.2) node[left] {$a$}--(3.3,1.2);
			\draw[mauChuDao,line width=0.9pt]
				(-1.5,-1) node[left] {$b$}--(3.3,-1);
			\draw[mauChuDao,line width=0.9pt]
				($(M)+(245:0.65)$)--($(N)+(65:0.65)$)
				node[above] {$c$};

			% Hai giao điểm
			\fill[red!75!black] (M) circle (1.5pt);
			\fill[red!75!black] (N) circle (1.5pt);
			\node[below right] at (M) {$M$};
			\node[above left] at (N) {$N$};

			% Các góc M_1, M_2, N_1, N_2
			\node at ($(M)+(145:0.43)$) {$1$};
			\node at ($(M)+(35:0.43)$) {$2$};
			\node at ($(N)+(215:0.43)$) {$2$};
			\node at ($(N)+(325:0.43)$) {$1$};
		\end{tikzpicture}
	\end{center}
\end{bt}

\begin{bt}
	Cho hình vẽ dưới đây, biết $a\parallel b$. Tính số đo $x$ và $y$.

	\begin{center}
		\begin{tikzpicture}[
			scale=1.15,
			line cap=round,
			line join=round,
			font=\small,
			every node/.style={inner sep=1.3pt}
		]
			\coordinate (A) at (0,1.2);
			\coordinate (B) at (2.0,1.2);
			\coordinate (C) at (-0.39,-1.0);
			\coordinate (D) at (3.27,-1.0);

			% Hai đường thẳng song song a và b
			\draw[mauChuDao,line width=0.9pt]
				(-1.1,1.2) node[left] {$a$}--(4.15,1.2);
			\draw[mauChuDao,line width=0.9pt]
				(-1.1,-1.0) node[left] {$b$}--(4.15,-1.0);

			% Hai cát tuyến c và d
			\draw[mauChuDao,line width=0.9pt]
				($(C)+(260:0.55)$)--($(A)+(80:0.65)$);
			\draw[mauChuDao,line width=0.9pt]
				($(B)+(120:0.65)$)--($(D)+(300:0.55)$);

			% Các giao điểm
			\foreach \P in {A,B,C,D}{
				\fill[red!75!black] (\P) circle (1.5pt);
			}

			% Tên điểm và tên hai cát tuyến
			\node[above left] at (A) {$A$};
			\node[above right] at (B) {$B$};
			\node[below right] at (C) {$C$};
			\node[below left] at (D) {$D$};
			\node[above] at ($(A)+(80:0.65)$) {$c$};
			\node[above] at ($(B)+(120:0.65)$) {$d$};

			% Số đo góc và các góc cần tìm
			\node[below right=4pt] at (A) {$100^\circ$};
			\node[below left=5pt] at (B) {$120^\circ$};
			\node[above right=4pt] at (C) {$x$};
			\node[below right=4pt] at (D) {$y$};
		\end{tikzpicture}
	\end{center}
\end{bt}

\subsection{Dạng 2. Chứng minh song song và tính góc}
\begin{bt}
	Cho hình vẽ dưới đây. Hai đường thẳng $PQ$ và $NO$ có song song với nhau
	không? Tại sao?


    \begin{center}
        \begin{tikzpicture}[
            scale=1.05,
            line cap=round,
            line join=round,
            font=\small,
            every node/.style={inner sep=1.3pt}
        ]
            % Các tọa độ
            \coordinate (O) at (0,0);
            \coordinate (N) at (-2.5,0);
            \coordinate (X) at (2.6,0);
            \coordinate (P) at ($(O)+(60:1.65)$);
            \coordinate (Q) at ($(P)+(3.1,0)$);

            % Sửa hướng OM thành 310 độ
            \coordinate (M) at ($(O)+(310:1.85)$);

            % Hai vùng tô xám
            \fill[gray!45,opacity=0.55] (P) circle (0.34);
            \fill[gray!45,opacity=0.55] (O) circle (0.42);

            % Các đường và tia
            \draw[mauChuDao,line width=0.9pt] (P)--(Q);
            \draw[mauChuDao,line width=0.9pt] (N)--(X);
            \draw[mauChuDao,line width=0.9pt] (P)--(O)--(M);

            % Góc OPQ = 120 độ
            \pic[
                draw=mauChuDao,
                "$120^\circ$",
                angle radius=5.5mm,
                angle eccentricity=1.45
            ] {angle=O--P--Q};

            % Góc POM = 110 độ
            \draw[mauChuDao]
                ($(O)+(60:0.31)$)
                arc[start angle=60,end angle=-50,radius=0.31];

            \node[fill=white,inner sep=1pt]
                at ($(O)+(12:0.85)$) {$110^\circ$};

            % Góc NOM = 130 độ
            \draw[mauChuDao]
                ($(O)+(180:0.38)$)
                arc[start angle=180,end angle=310,radius=0.38];

            \draw[mauChuDao]
                ($(O)+(180:0.47)$)
                arc[start angle=180,end angle=310,radius=0.47];

            \node at ($(O)+(242:0.82)$) {$130^\circ$};

            % Các điểm và nhãn
            \foreach \Pnt in {N,P,O,Q,M}{
                \fill[red!75!black] (\Pnt) circle (1.35pt);
            }

            \node[above left=1pt] at (P) {$P$};
            \node[above left=1pt] at (O) {$O$};
            \node[left=2pt] at (N) {$N$};
            \node[right=3pt] at (Q) {$Q$};
            \node[right=2pt] at (X) {$x$};
            \node[below right=1pt] at (M) {$M$};

        \end{tikzpicture}
    \end{center}

\end{bt}
\begin{bt}
	Cho đoạn thẳng $AB$. Trên cùng một nửa mặt phẳng có bờ là đường thẳng
	$AB$, vẽ các tia $Ax$, $By$ sao cho
	$\widehat{ABy}=\alpha$ và $\widehat{xAB}=\dfrac{7}{2}\alpha$.
	Tính $\alpha$ để $Ax\parallel By$.
\end{bt}

\begin{bt}
	Theo hình vẽ, cho $\widehat{xOy}=70^\circ$ 
	$(\widehat{O_1}=70^\circ)$, $\widehat{A_1}=110^\circ$ và
	$\widehat{B_2}=110^\circ$. Hãy chứng minh:
	\begin{enumerate}[label=\alph*)]
		\item $a\parallel Ox$;
		\item $b\parallel Oy$.
	\end{enumerate}

    \begin{center}
        \begin{tikzpicture}[
            scale=1.05,
            line cap=round,
            line join=round,
            font=\small,
            every node/.style={inner sep=1.3pt}
        ]
            \coordinate (O) at (0,0);
            \coordinate (X) at (3.1,0);
            \coordinate (A) at ($(O)+(70:2.25)$);
            \coordinate (B) at ($(A)+(-2.35,0)$);

            % Tạo điểm Y trên tia Oy (kéo dài hơn A)
            \coordinate (Y) at ($(O)!1.48!(A)$);

            % Hai điểm phụ để vẽ đường thẳng b song song tuyệt đối với Oy
% Hai điểm gốc theo phương Oy
            \coordinate (Bleft)  at ($(B)+(O)-(Y)$);
            \coordinate (Bright) at ($(B)+(Y)-(O)$);

            % Lấy gần B hơn để đường b ngắn lại
            \coordinate (B1) at ($(B)!0.58!(Bleft)$);
            \coordinate (B2) at ($(B)!0.58!(Bright)$);

            % Đường thẳng a và tia Ox
            \draw[mauChuDao,line width=0.9pt]
                ($(B)+(-1.1,0)$)--($(A)+(1.35,0)$)
                node[above left=1pt] {$a$};

            \draw[mauChuDao,line width=0.9pt] (O)--(X)
                node[below left=1pt] {$x$};

            % Tia Oy đi qua A
            \draw[mauChuDao,line width=0.9pt] (O)--(Y);
            \node[above right] at (Y) {$y$};

            % Đường thẳng b đi qua B và song song với Oy
            \draw[mauChuDao,line width=0.9pt] (B1)--(B2);
            \node[above right] at ($(B)!0.98!(B2)$) {$b$};

            % Các giao điểm
            \foreach \Pnt in {O,A,B}{
                \fill[red!75!black] (\Pnt) circle (1.45pt);
            }

            \node[below left=1pt] at (O) {$O$};
            \node[above left=1pt] at (A) {$A$};
            \node[above right=2pt] at (B) {$B$};

            % Ký hiệu các góc O_1, A_1 và B_2
            \node at ($(O)+(35:0.28)$) {$1$};
            \node at ($(A)+(-55:0.28)$) {$1$};
            \node at ($(B)+(130:0.28)$) {$2$};
        \end{tikzpicture}
    \end{center}

\end{bt}
\subsection{Dạng 3. Tổng hợp tính góc và chứng minh song song}
\begin{bt}
	Cho hình vẽ. Tìm trên hình các đường thẳng song song với $OC$.
	Vì sao?

	\begin{center}
		\begin{tikzpicture}[
			scale=1.05,
			line cap=round,
			line join=round,
			font=\small,
			every node/.style={inner sep=1.3pt}
		]
			\coordinate (O) at (0,0);
			\coordinate (C) at (-2.05,0);
			\coordinate (D) at (0,-1.55);
			\coordinate (E) at (2.35,-1.55);
			\coordinate (A) at ($(O)+(50:1.35)$);
			\coordinate (B) at ($(A)+(2.35,0)$);

			% Vùng tô xám quanh hai đỉnh có góc đã cho
			\fill[gray!45,opacity=0.55] (O) circle (0.36);
			\fill[gray!45,opacity=0.55] (A) circle (0.36);

			% Các đường và đoạn thẳng
			\draw[mauChuDao,line width=0.9pt] (C)--(O)--(A)--(B);
			\draw[mauChuDao,line width=0.9pt] (O)--(D)--(E);

			% Góc 140 độ tại O
			\draw[mauChuDao] ($(O)+(50:0.28)$)
				arc[start angle=50,end angle=-90,radius=0.28];
			\node[right=3pt] at ($(O)+(-23:0.48)$) {$140^\circ$};

			% Góc 130 độ tại A với các cung đồng tâm
			\draw[mauChuDao] ($(A)+(230:0.27)$)
				arc[start angle=230,end angle=360,radius=0.27];
			\draw[mauChuDao] ($(A)+(230:0.34)$)
				arc[start angle=230,end angle=360,radius=0.34];
			\draw[mauChuDao] ($(A)+(230:0.41)$)
				arc[start angle=230,end angle=360,radius=0.41];
			\node[below right=5pt] at (A) {$130^\circ$};

			% Dấu vuông: OD vuông góc OC và OD vuông góc DE
			\draw[mauChuDao,line width=0.8pt]
				($(O)+(-0.18,0)$)--($(O)+(-0.18,-0.18)$)--($(O)+(0,-0.18)$);
			\draw[mauChuDao,line width=0.8pt]
				($(D)+(0,0.18)$)--($(D)+(0.18,0.18)$)--($(D)+(0.18,0)$);

			% Các điểm và nhãn
			\foreach \Pnt in {C,O,A,B,D,E}{
				\fill[red!75!black] (\Pnt) circle (1.35pt);
			}
			\node[left=3pt] at (C) {$C$};
			\node[above left=1pt] at (O) {$O$};
			\node[above left=1pt] at (A) {$A$};
			\node[right=3pt] at (B) {$B$};
			\node[below left=1pt] at (D) {$D$};
			\node[right=3pt] at (E) {$E$};
		\end{tikzpicture}
	\end{center}
\end{bt}

\begin{bt}
	Cho hình vẽ sau. Biết
	$\widehat{AMN}=2\widehat{BMN}$ và
	$\widehat{MND}-\widehat{MNC}=60^\circ$.
	Chứng minh rằng $AB\parallel CD$.

	\begin{center}
		\begin{tikzpicture}[
			scale=1.1,
			line cap=round,
			line join=round,
			font=\small,
			every node/.style={inner sep=1.3pt}
		]
			\coordinate (M) at (0,1.25);
			\coordinate (N) at (1.65,-0.65);
			\coordinate (A) at (-1.45,1.25);
			\coordinate (B) at (3.0,1.25);
			\coordinate (C) at (-1.65,-0.65);
			\coordinate (D) at (3.15,-0.65);

			% Hai đường thẳng AB, CD và cát tuyến MN
			\draw[mauChuDao,line width=0.9pt] (A)--(B);
			\draw[mauChuDao,line width=0.9pt] (C)--(D);
			\draw[mauChuDao,line width=0.9pt] (M)--(N);

			% Hai giao điểm
			\fill[red!75!black] (M) circle (1.5pt);
			\fill[red!75!black] (N) circle (1.5pt);

			% Nhãn các điểm
			\node[above left=1pt] at (A) {$A$};
			\node[above=3pt] at (M) {$M$};
			\node[above left=1pt] at (B) {$B$};
			\node[above left=1pt] at (C) {$C$};
			\node[above right=1pt] at (N) {$N$};
			\node[above left=1pt] at (D) {$D$};
		\end{tikzpicture}
	\end{center}
\end{bt}

\begin{bt}
	Cho hình vẽ sau. Chứng minh rằng $AB\parallel CD$.

	\begin{center}
		\begin{tikzpicture}[
			scale=1.1,
			line cap=round,
			line join=round,
			font=\small,
			every node/.style={inner sep=1.3pt}
		]
			\coordinate (A) at (0,0);
			\coordinate (B) at (3.1,0);
			\coordinate (P) at ($(A)+(120:1.65)$);
			\coordinate (C) at ($(A)+(260:1.55)$);
			\coordinate (D) at ($(C)+(3.35,0)$);

			% Các đoạn thẳng trong hình
			\draw[mauChuDao,line width=0.9pt] (P)--(A)--(B);
			\draw[mauChuDao,line width=0.9pt] (A)--(C)--(D);

			% Các góc đã cho tại A và C
			\pic[draw=mauChuDao,"$120^\circ$",angle radius=2.5mm,
				angle eccentricity=2.5] {angle=B--A--P};
			\pic[draw=mauChuDao,"$140^\circ$",angle radius=3.5mm,
				angle eccentricity=2] {angle=P--A--C};
			\pic[draw=mauChuDao,"$80^\circ$",angle radius=4.5mm,
				angle eccentricity=1.7] {angle=D--C--A};

			% Các điểm
			\foreach \Pnt in {P,A,B,C,D}{
				\fill[red!75!black] (\Pnt) circle (1.5pt);
			}

			% Nhãn các điểm
			\node[below right=2pt] at (A) {$A$};
			\node[above left=1pt] at (B) {$B$};
			\node[below left=1pt] at (C) {$C$};
			\node[below left=1pt] at (D) {$D$};
		\end{tikzpicture}
	\end{center}
\end{bt}
